\section{Формулирование требований к программному комплексу}

Проектирование любого программного продукта начинается с определения набора требований, которые описывают, что система должна делать и какими свойствами обладать. Требования делятся на две основные категории: функциональные и нефункциональные.

\subsection{Функциональные требования}

Функциональные требования определяют конкретные функции, которые система должна предоставлять пользователю. На основе анализа предметной области были сформулированы следующие требования:

\begin{enumerate}
    \item \textbf{Мониторинг системы в реальном времени.} Система должна непрерывно собирать и отображать следующие метрики:
    \begin{itemize}
        \item Общую и по-ядерную загрузку центрального процессора.
        \item Информацию о модели ЦП, его частоте и количестве ядер.
        \item Объем использованной и доступной оперативной памяти.
        \item Текущую скорость операций чтения и записи на дисковых накопителях.
        \item Текущую скорость приема и передачи данных по сети.
        \item Общую информацию об операционной системе и времени ее работы (uptime).
    \end{itemize}
    
    \item \textbf{Проведение синтетических тестов ЦП.} Система должна предоставлять возможность запускать тесты для оценки производительности процессора, включая:
    \begin{itemize}
        \item Тест на скорость выполнения целочисленных операций.
        \item Тест на скорость выполнения операций с плавающей запятой.
    \end{itemize}
    
    \item \textbf{Проведение тестов подсистемы памяти.} Система должна включать бенчмарки для оценки различных аспектов иерархии памяти:
    \begin{itemize}
        \item Тесты для измерения задержки (latency) доступа к кэшам L1, L2, L3 и к основной оперативной памяти.
        \item Тесты для измерения пропускной способности (bandwidth) оперативной памяти при операциях чтения и записи.
        \item Тесты для измерения производительности постоянных накопителей (SSD/HDD).
    \end{itemize}

    \item \textbf{Проведение тестов механизмов параллелизма.} Система должна оценивать эффективность встроенных в язык программирования средств для параллельной обработки данных.
    
    \item \textbf{Графический пользовательский интерфейс.} Все функции мониторинга и тестирования должны быть доступны через единый, интуитивно понятный графический интерфейс. Интерфейс должен наглядно представлять данные и обеспечивать удобное управление.
\end{enumerate}

\subsection{Нефункциональные требования}

Нефункциональные требования определяют атрибуты качества и ограничения для системы. Они не менее важны, чем функциональные, так как определяют, насколько хорошо система выполняет свои функции.

\begin{itemize}
    \item \textbf{Кросс-платформенность.} Приложение должно корректно работать на основных операционных системах: Windows, macOS и Linux.
    \item \textbf{Низкие накладные расходы.} Программный комплекс, предназначенный для анализа производительности, сам должен оказывать минимальное влияние на измеряемую систему.
    \item \textbf{Отзывчивость интерфейса.} Графический интерфейс должен оставаться отзывчивым и не «зависать» во время выполнения длительных тестов или сбора данных.
    \item \textbf{Простота развертывания.} Приложение должно распространяться в виде единого исполняемого файла, не требующего от пользователя установки дополнительных зависимостей или библиотек.
    \item \textbf{Модульность и расширяемость.} Архитектура должна позволять легко добавлять новые тесты или модули мониторинга в будущем.
\end{itemize}

\section{Обоснование выбора технологического стека}

На основе сформулированных требований был проведен анализ и выбор технологического стека. Ключевыми компонентами стека являются язык программирования и библиотека для создания графического интерфейса.

\subsection{Выбор языка программирования}

Рассматривались следующие кандидаты: C++, Python, Rust и Go.

\begin{itemize}
    \item \textbf{C++:} Традиционный выбор для высокопроизводительных системных утилит. Обеспечивает максимальную производительность и полный контроль над ресурсами. Однако его использование сопряжено с высокой сложностью разработки, ручным управлением памятью и проблемами при обеспечении кросс-платформенной сборки GUI.
    \item \textbf{Python:} Популярен для быстрой разработки, имеет отличные библиотеки для анализа данных. Главный недостаток — низкая производительность интерпретатора и наличие глобальной блокировки интерпретатора (GIL), что затрудняет написание по-настоящему параллельного кода для утилизации всех ядер ЦП. Это делает его непригодным для написания точных бенчмарков.
    \item \textbf{Rust:} Современный язык, ориентированный на безопасность и производительность, сопоставимую с C++. Обладает сложной моделью владения и заимствования, что значительно повышает порог вхождения и замедляет разработку по сравнению с Go.
    \item \textbf{Go (Golang):} Выбран как оптимальный компромисс. Он предлагает производительность, близкую к C++ и Rust, но при этом значительно проще в освоении. Его модель параллелизма с горутинами и каналами идеально подходит для задач асинхронного обновления GUI и фонового выполнения тестов. Встроенная кросс-компиляция и статическая линковка полностью удовлетворяют требованиям к простоте развертывания.
\end{itemize}

\subsection{Выбор библиотеки для GUI}

Для языка Go существует несколько подходов к созданию GUI. Рассматривались Wails, Gio UI и Fyne.

\begin{itemize}
    \item \textbf{Wails:} Позволяет использовать веб-технологии (HTML/CSS/JS) для создания интерфейса. Это дает большую гибкость в дизайне, но требует знаний веб-стека и может приводить к большему потреблению памяти и оверхеду из-за встроенного веб-движка.
    \item \textbf{Gio UI:} Низкоуровневая библиотека, использующая модель немедленного рендеринга (immediate mode). Предоставляет максимальную производительность и гибкость, но требует от разработчика написания большего количества кода для управления состоянием и отрисовкой.
    \item \textbf{Fyne:} Выбран как наиболее сбалансированное решение. Это чисто нативный инструментарий, написанный на Go, который не использует веб-технологии. Он предлагает декларативный подход (retained mode), богатый набор готовых виджетов и автоматическое соответствие внешнему виду целевой ОС. Это позволило быстро разработать функциональный и современно выглядящий интерфейс, полностью удовлетворив требования проекта.
\end{itemize}

\section{Проектирование архитектуры}

На основе требований и выбранного стека была спроектирована многоуровневая модульная архитектура.

\subsection{Высокоуровневая структура}

Приложение разделено на три основных слоя:
\begin{enumerate}
    \item \textbf{Слой представления (GUI Layer):} Отвечает за отображение информации и взаимодействие с пользователем. Реализован в пакете \texttt{pkg/gui}.
    \item \textbf{Слой бизнес-логики (Business Logic Layer):} Содержит логику выполнения бенчмарков и сбора системных метрик. Реализован в пакетах \texttt{pkg/benchmark}, \texttt{pkg/memory}, \texttt{pkg/concurrency} и \texttt{pkg/profiling}.
    \item \textbf{Слой доступа к данным (Data Access Layer):} В данном проекте этот слой представлен библиотекой \texttt{gopsutil}, которая абстрагирует системные вызовы ОС.
\end{enumerate}

Такое разделение позволяет каждому слою быть независимым. Например, тесты из слоя бизнес-логики можно запускать отдельно, без графического интерфейса, что упрощает их отладку.

\subsection{Проектирование основных модулей}

\paragraph{Модуль GUI.} Спроектирован на основе вкладочной навигации (\texttt{container.NewAppTabs}). Каждая вкладка представляет собой отдельную панель, создаваемую собственной функцией-конструктором (например, \texttt{NewDashboardPanel}). Внутри панелей информация группируется с помощью виджетов-карточек (\texttt{widget.NewCard}) для логического разделения. Все длительные операции запускаются в отдельных горутинах, чтобы не блокировать поток GUI, что является ключевым архитектурным решением для обеспечения отзывчивости.

\paragraph{Модуль запуска тестов.} Для унификации процесса тестирования был спроектирован универсальный обработчик \texttt{benchmark.TestRunner}. Он принимает на вход любую функцию, удовлетворяющую интерфейсу \texttt{BenchmarkFunc}, и управляет ее многократным запуском, сбором статистики и передачей прогресса через канал. Это позволило избежать дублирования кода и отделить логику проведения теста от логики самого теста.
